% Comercial I - edificio corporativo com nucleo, acessos e SPDA
\section{Projeto Comercial I — edifício corporativo com SPDA}

\newcommand{\EdifComBase}{%
\draw[arqExt] (0,0) rectangle (26,18);
\ArqParedeInt{0}{8}{26}{8}\ArqParedeInt{0}{10}{26}{10}
\foreach \x in {5.2,10.4,15.6,20.8}{\ArqParedeInt{\x}{10}{\x}{18}}
\ArqParedeInt{5}{0}{5}{8}\ArqParedeInt{10}{0}{10}{8}\ArqParedeInt{18}{0}{18}{8}\ArqParedeInt{22}{0}{22}{8}
% nucleo: escada, elevador, shaft e sanitarios com corredor lateral de 1,30 m
\ArqParedeInt{13.3}{0}{13.3}{8}\ArqParedeInt{16.7}{0}{16.7}{8}
\ArqParedeInt{14.3}{0}{14.3}{5.8}
\ArqParedeInt{14.3}{1.8}{16.7}{1.8}\ArqParedeInt{14.3}{3.7}{16.7}{3.7}\ArqParedeInt{14.3}{5.8}{16.7}{5.8}
% portas salas superiores e inferiores
\foreach \x in {1.0,6.2,11.4,16.6,21.8}{\ArqPortaHUp{\x}{10}{0.90}}
\foreach \x in {1.0,6.0,18.7,22.7}{\ArqPortaHDown{\x}{8}{0.90}}
% escada, elevador e shaft acessados diretamente pelo corredor central
\ArqPortaHDown{11.2}{8}{0.90}\ArqPortaCorrerH{13.75}{15.25}{8}\ArqPortaHDown{15.75}{8}{0.70}
% corredor dos sanitarios: acesso independente do patamar da escada
\ArqVaoH{16.85}{17.95}{8}
\ArqPortaVLeft{16.7}{0.40}{0.75}\ArqPortaVLeft{16.7}{2.25}{0.75}\ArqPortaVLeft{16.7}{4.25}{0.90}
% fachada
\foreach \xa/\xb in {0.8/4.3,5.9/9.5,11.1/14.7,16.3/19.9,21.5/25.2}{\ArqJanH{\xa}{\xb}{18}}
\foreach \xa/\xb in {0.8/4.2,5.8/9.2,18.5/21.4,22.5/25.2}{\ArqJanH{\xa}{\xb}{0}}
% mobiliario salas superiores
\foreach \x in {0.6,5.8,11.0,16.2,21.4}{\ArqDesk{\x}{12.0}{1.45}{0.60}\ArqDesk{\x+2.3}{12.0}{1.45}{0.60}\ArqDesk{\x}{15.2}{1.45}{0.60}\ArqDesk{\x+2.3}{15.2}{1.45}{0.60}}
% salas inferiores
\foreach \x in {0.5,5.5,18.4}{\ArqDesk{\x}{2.0}{1.25}{0.58}\ArqDesk{\x+2.0}{2.0}{1.25}{0.58}\ArqDesk{\x}{5.2}{1.25}{0.58}\ArqDesk{\x+2.0}{5.2}{1.25}{0.58}}
\ArqMesa{22.5}{2.4}{2.8}{1.35}\ArqBancada{22.5}{6.6}{2.8}{0.45}
% nucleo vertical e sanitarios
\ArqEscada{10.25}{0.35}{2.75}{7.25}
\ArqElevador{13.65}{6.05}{1.75}{1.70}
\draw[arqEquip] (15.60,6.05) rectangle (16.55,7.75);\node[arqSub,align=center] at (16.08,6.90){SHAFT\\TÉCN.};
\ArqVaso{15.45}{0.65}\ArqPia{14.55}{1.18}{0.60}{0.35}\node[arqSub] at (15.45,1.55){WC-M};
\ArqVaso{15.45}{2.55}\ArqPia{14.55}{3.05}{0.60}{0.35}\node[arqSub] at (15.45,3.42){WC-F};
\ArqVaso{15.50}{4.55}\ArqPia{14.55}{5.15}{0.70}{0.38}\node[arqSub] at (15.45,5.55){WC-PCD};
\node[arqSub,rotate=90] at (17.35,3.0){ACESSO WC};
\node[arqSub] at (14.45,7.55){HALL ELEV.};
% rotulos
\foreach \x/\n in {2.6/01,7.6/02,13.0/03,18.2/04,23.4/05}{\node[arqLabel] at (\x,17.35){SALA \n};}
\node[arqLabel] at (2.5,0.65){SALA 06};\node[arqLabel] at (7.5,0.65){SALA 07};\node[arqLabel] at (20.0,0.65){SALA 08};\node[arqLabel] at (24.0,0.65){REUNIÃO};
\node[arqSub] at (5.0,9.0){CORREDOR - 2,00 m};\node[arqSub] at (20.5,9.0){CORREDOR};
\CotaH{0}{26}{18.22}{26,00 m}\CotaV{0}{18}{26.22}{18,00 m}
}

\begin{frame}{Comercial I — pavimento típico refinado}
\centering\begin{tikzpicture}[scale=0.37,every node/.style={font=\tiny}]\EdifComBase\end{tikzpicture}
\vspace{0.4mm}\scriptsize Salas voltadas ao corredor central; escada, elevador e shaft técnico com acessos próprios; sanitários atendidos por corredor lateral de 1,30 m. Potência e telecom usam compartimentação/segregação compatível no shaft.
\end{frame}

\begin{frame}{Comercial I — organização dos quatro pavimentos}
\small\begin{columns}[T,onlytextwidth]
\column{0.24\textwidth}\begin{caixaProjeto}[title=Térreo]\begin{itemize}\item lobby;\item 6 salas;\item reunião;\item WC M/F/PCD;\item elevador + escada.\end{itemize}\end{caixaProjeto}
\column{0.24\textwidth}\begin{caixaProjeto}[title=1º pav.]\begin{itemize}\item 8 salas;\item reunião;\item WC M/F/PCD;\item shaft/IT.\end{itemize}\end{caixaProjeto}
\column{0.24\textwidth}\begin{caixaProjeto}[title=2º pav.]\begin{itemize}\item 8 salas;\item reunião;\item WC M/F/PCD;\item shaft/IT.\end{itemize}\end{caixaProjeto}
\column{0.24\textwidth}\begin{caixaProjeto}[title=3º pav.]\begin{itemize}\item 8 salas;\item reunião;\item WC M/F/PCD;\item acesso técnico.\end{itemize}\end{caixaProjeto}
\end{columns}\end{frame}

\begin{frame}{Comercial I — elétrica do pavimento típico}
\centering\begin{tikzpicture}[scale=0.37,every node/.style={font=\tiny}]
\EdifComBase\SimQD{16.05}{9.65}{QDP}
\foreach \x in {1.3,3.9,6.5,9.1,11.7,14.3,16.9,19.5,22.1,24.7}{\SimLuz{\x}{14.0}}
\foreach \x in {1.3,3.9,6.5,9.1,18.7,20.7,22.7,24.7}{\SimLuz{\x}{4.4}}
\foreach \x in {1.5,4.5,7.5,10.5,13.5,16.5,19.5,22.5,25.0}{\SimLuz{\x}{9.0}}
\foreach \x/\y in {0.4/11,4.8/11,5.6/11,9.8/11,10.8/11,15/11,16/11,20.2/11,21.2/11,25.6/11,0.4/7.2,4.6/7.2,5.4/7.2,9.6/7.2,18.4/7.2,21.5/7.2,22.3/7.2,25.6/7.2}{\SimTUG{\x}{\y}\SimDados{\x}{\y+0.35}}
\foreach \x in {2.6,7.6,13,18.2,23.4}{\SimTUE{\x}{17.6}{AC}}\foreach \x in {2.5,7.5,20,24}{\SimTUE{\x}{0.4}{AC}}\SimTUE{14.5}{7.5}{EL}
\foreach \x in {2,7,12,17,22}{\SimEmerg{\x}{9.0}}\SimEmerg{11.6}{7.55}\SimEmerg{14.5}{7.55}
\draw[corPrim!70,dashed,thick] (16.05,9.65)--(16.05,9)--(2,9);\draw[corPrim!70,dashed,thick] (16.05,9)--(25,9);
\end{tikzpicture}\end{frame}

\begin{frame}[shrink=5]{Comercial I — carga e demanda do estudo}
\scriptsize
\begin{tabularx}{\textwidth}{>{\bfseries}p{3.0cm}c c c c X}
\toprule
Grupo & Dado instalado & $fp$ & $P_{inst}$ & $f_d$ & $P_{dem}$\\
\midrule
Iluminação/emergência & 18,0 kW & 0,95 & 18,0 kW & 0,90 & 16,2 kW\\
TUG/TI & 32,0 kVA & 0,90 & 28,8 kW & 0,70 & 20,2 kW\\
Climatização & 72,0 kW & 0,92 & 72,0 kW & 0,80 & 57,6 kW\\
Elevador & 18,0 kW & 0,85 & 18,0 kW & 0,80 & 14,4 kW\\
Serviços comuns & 20,0 kW & 0,90 & 20,0 kW & 0,60 & 12,0 kW\\
\midrule
\textbf{Total} & -- & -- & \textbf{156,8 kW} & -- & \textbf{120,4 kW}\\
\bottomrule
\end{tabularx}
\begin{caixaAt}
Os fatores de demanda e de potência são \textbf{premissas didáticas}, não uma tabela universal. No executivo, usar ocupação, simultaneidade, dados de placa, estratégia do elevador/HVAC e método aceito pela Equatorial.
\end{caixaAt}
\end{frame}

\begin{frame}{Comercial I — distribuição vertical e conexão}
\small\begin{columns}[T,onlytextwidth]
\column{0.49\textwidth}
\begin{caixaProjeto}[title=Ordem de grandeza]
Pelas premissas do estudo:
\[
S_{dem}\approx\SI{132}{kVA},\qquad
I_{dem}\approx\frac{132000}{\sqrt3\,380}=\SI{201}{A}.
\]
QDI, QDT/TI, QDAC, QD-ELEV e QD-SERV partem do QGBT com alimentadores, proteção e seletividade próprios.
\end{caixaProjeto}
\column{0.49\textwidth}
\begin{caixaNorma}[title=Interface correta]
O empreendimento não pode ser enquadrado pela tabela de unidade individual da NT.00001. Medição coletiva, demanda do conjunto e prumadas seguem a \textbf{NT.00004.EQTL Rev.08}; se a solução de conexão for em MT, aplicar também a \textbf{NT.00002.EQTL Rev.10} e os dados fornecidos pela distribuidora.
\end{caixaNorma}
\end{columns}\end{frame}

\begin{frame}{Comercial I — cobertura e SPDA}
\centering\begin{tikzpicture}[scale=0.36,every node/.style={font=\tiny}]
\draw[arqExt] (0,0) rectangle (26,18);\draw[arqEquip,line width=0.8pt] (10.2,5.2) rectangle (17.2,10.8);\node[arqLabel] at (13.7,8){CASA DE MÁQUINAS\\ACESSO À ESCADA};
\draw[arqEquip] (19,12) rectangle (23.5,15);\node[arqSub] at (21.25,13.5){CONDENSADORAS};
\foreach \y in {0,6,12,18}{\draw[corAt!80,thick] (0,\y)--(26,\y);}\foreach \x in {0,6.5,13,19.5,26}{\draw[corAt!80,thick] (\x,0)--(\x,18);}
\foreach \x/\y in {0/0,6.5/0,13/0,19.5/0,26/0,26/6,26/12,26/18,19.5/18,13/18,6.5/18,0/18,0/12,0/6}{\fill[corNorma] (\x,\y) circle (0.14);}
\end{tikzpicture}\vspace{0.3mm}\scriptsize Malha e descidas são conceituais; classe e espaçamentos finais decorrem da análise de risco e da série ABNT NBR 5419 aplicável.
\end{frame}

\begin{frame}{Comercial I — SPDA, DPS e aterramento integrados}
\centering\begin{tikzpicture}[scale=0.70,every node/.style={font=\tiny}]
\tikzset{b/.style={draw=corPrim,fill=corDef!7,rounded corners,minimum width=2.2cm,minimum height=0.58cm,align=center}}
\node[b](r)at(-4,3.5){Rede Equatorial};\node[b](qg)at(-4,2.3){QGBT\\DPS origem};\node[b](qd)at(-4,1.1){QDs pavimentos};\node[b](sp)at(1,3.5){SPDA cobertura};\node[b](desc)at(1,2.3){Descidas};\node[b](bep)at(1,1.1){BEP / equipot.};\node[b](malha)at(1,-0.1){aterramento};
\draw[-{Stealth},thick,corPrim](r)--(qg)--(qd);\draw[-{Stealth},thick,corAt](sp)--(desc)--(bep)--(malha);\draw[corNorma,thick](qg)--(bep);\draw[corNorma,thick](qd)--(bep);
\end{tikzpicture}\end{frame}
